\chapter{无穷级数与无穷积分}


\section{无穷级数}

无穷级数和无穷积分分别是离散形式和连续形式的求和.判断它们的敛散性时,真正要控制的不是前面有限多项,而是尾部.对级数而言尾部是
\begin{align*}
	\sum_{n=N_1}^{N_2}u_n,
\end{align*}
对积分而言尾部是
\begin{align*}
	\int_{A_1}^{A_2}f(x)\,\dd x.
\end{align*}
因此,无论是逐点收敛、绝对收敛还是一致收敛,核心都可以看成同一个 Cauchy 思想在离散和连续两种情形下的表现.

\section{基本概念}

\textbf{不含参无穷级数的敛散性.}
设 $\{a_n\}$ 为数列.若部分和
\begin{align*}
	S_N=\sum_{n=1}^N a_n
\end{align*}
当 $N\to\infty$ 时存在有限极限,则称无穷级数
\begin{align*}
	\sum_{n=1}^{\infty}a_n
\end{align*}
收敛;否则称其发散.等价地,对任意 $\varepsilon>0$,存在 $N_0$,使得当 $N_2>N_1\geq N_0$ 时,
\begin{align*}
	\left|\sum_{n=N_1}^{N_2}a_n\right|<\varepsilon.
\end{align*}
这就是级数的 Cauchy 收敛准则.

若
\begin{align*}
	\sum_{n=1}^{\infty}|a_n|
\end{align*}
收敛,则称原级数绝对收敛;若原级数收敛而绝对值级数发散,则称原级数条件收敛.

\textbf{含参无穷级数的敛散性.}
设 $E$ 为参数集合,且 $u_n(t)$ 在 $E$ 上定义.若对每一个固定的 $t\in E$,普通级数
\begin{align*}
	\sum_{n=1}^{\infty}u_n(t)
\end{align*}
都收敛,则称该含参级数在 $E$ 上逐点收敛.

进一步,若对任意 $\varepsilon>0$,存在 $N_0$,使得当 $N_2>N_1\geq N_0$ 时,对一切 $t\in E$ 均有
\begin{align*}
	\left|\sum_{n=N_1}^{N_2}u_n(t)\right|<\varepsilon,
\end{align*}
则称该含参级数在 $E$ 上一致收敛.这就是含参级数的 Cauchy 一致收敛准则.

若
\begin{align*}
	\sum_{n=1}^{\infty}|u_n(t)|
\end{align*}
在 $E$ 上一致收敛,则称原级数在 $E$ 上一致绝对收敛.若原级数在 $E$ 上一致收敛,但绝对值级数在 $E$ 上不一致收敛,则称原级数在 $E$ 上一致条件收敛.

\section{常用判别法}

\textbf{比较判别法.}
设从某一项起 $0\leq a_n\leq b_n$.若
\begin{align*}
	\sum_{n=1}^{\infty}b_n
\end{align*}
收敛,则
\begin{align*}
	\sum_{n=1}^{\infty}a_n
\end{align*}
也收敛;若
\begin{align*}
	\sum_{n=1}^{\infty}a_n
\end{align*}
发散,则
\begin{align*}
	\sum_{n=1}^{\infty}b_n
\end{align*}
也发散.

极限形式如下.设从某一项起 $a_n\geq0$, $b_n>0$.若
\begin{align*}
	\lim_{n\to\infty}\frac{a_n}{b_n}=L,
\end{align*}
则当 $0<L<\infty$ 时,两级数同时收敛或同时发散.

含参形式如下.若对一切 $n$ 和一切 $t\in E$ 有
\begin{align*}
	|u_n(t)|\leq M_n,
\end{align*}
且
\begin{align*}
	\sum_{n=1}^{\infty}M_n
\end{align*}
收敛,则
\begin{align*}
	\sum_{n=1}^{\infty}u_n(t)
\end{align*}
在 $E$ 上一致绝对收敛.这也称为 Weierstrass 判别法.

\textbf{Dirichlet 判别法.}
设 $a_n,b_n$ 为数列.若存在常数 $M>0$,使得对一切 $N$ 都有
\begin{align*}
	\left|\sum_{n=1}^N a_n\right|\leq M,
\end{align*}
又若 $b_n$ 单调且
\begin{align*}
	\lim_{n\to\infty}b_n=0,
\end{align*}
则
\begin{align*}
	\sum_{n=1}^{\infty}a_nb_n
\end{align*}
收敛.

含参形式如下.若存在常数 $M>0$,使得对一切 $N$ 和一切 $t\in E$ 都有
\begin{align*}
	\left|\sum_{n=1}^N a_n(t)\right|\leq M,
\end{align*}
并且 $b_n(t)$ 对每个固定的 $t$ 关于 $n$ 单调,且
\begin{align*}
	\lim_{n\to\infty}\sup_{t\in E}|b_n(t)|=0,
\end{align*}
则
\begin{align*}
	\sum_{n=1}^{\infty}a_n(t)b_n(t)
\end{align*}
在 $E$ 上一致收敛.

\textbf{Abel 判别法.}
若
\begin{align*}
	\sum_{n=1}^{\infty}a_n
\end{align*}
收敛,且 $b_n$ 单调有界,则
\begin{align*}
	\sum_{n=1}^{\infty}a_nb_n
\end{align*}
收敛.

含参形式如下.若
\begin{align*}
	\sum_{n=1}^{\infty}a_n(t)
\end{align*}
在 $E$ 上一致收敛,且 $b_n(t)$ 对每个固定的 $t$ 关于 $n$ 单调,并且存在常数 $M>0$,使得
\begin{align*}
	|b_n(t)|\leq M,
	\qquad n\geq1,
	\quad t\in E,
\end{align*}
则
\begin{align*}
	\sum_{n=1}^{\infty}a_n(t)b_n(t)
\end{align*}
在 $E$ 上一致收敛.

\textbf{常用幂级数尾部.}
幂函数在级数中对应 $p$ 级数:
\begin{align*}
	\sum_{n=1}^{\infty}\frac1{n^p}
\end{align*}
在 $p>1$ 时收敛,在 $p\leq1$ 时发散.若 $f(x)$ 在 $[1,\infty)$ 上非负单调递减,则
\begin{align*}
	\sum_{n=1}^{\infty}f(n)
\end{align*}
与
\begin{align*}
	\int_1^\infty f(x)\,\dd x
\end{align*}
同时收敛或同时发散.这个积分判别法说明,在正项情形下,级数和无穷积分只是同一种尾部估计的两种写法.

\begin{exbox}
	考虑正项级数
	\begin{align*}
		\sum_{n=1}^{\infty}\frac1{n^p}.
	\end{align*}
	由于 $x^{-p}$ 在 $[1,\infty)$ 上非负单调递减,积分判别法给出
	\begin{align*}
		\sum_{n=1}^{\infty}\frac1{n^p}
		\quad\hbox{与}\quad
		\int_1^\infty x^{-p}\,\dd x
	\end{align*}
	同时收敛或同时发散.而
	\begin{align*}
		\int_1^\infty x^{-p}\,\dd x
	\end{align*}
	在 $p>1$ 时收敛,在 $p\leq1$ 时发散.因此
	\begin{gather*}
		\sum_{n=1}^{\infty}\frac1{n^p}
		\begin{cases}
			\text{收敛},&p>1,\\[3pt]
			\text{发散},&p\leq1.
		\end{cases}
	\end{gather*}
\end{exbox}

\textbf{下面考虑一种典型的级数形式,它是后面振荡积分的离散对应物.}

以下令
\begin{align*}
	S(t,p)=\sum_{n=1}^{\infty}\frac{e^{\mathrm{i}nt}}{n^p}.
\end{align*}
作为含参级数讨论时,取
\begin{align*}
	E\subset\mathbb R\setminus2\pi\mathbb Z,
\end{align*}
并记
\begin{align*}
	\delta_E=\inf_{t\in E}\operatorname{dist}(t,2\pi\mathbb Z).
\end{align*}
其中 $\delta_E=0$ 表示参数可以趋近于某个 $2\pi$ 的整数倍,这正是离散振荡和连续振荡中共同会破坏一致性的地方.

\begin{exbox}
	考虑振荡级数
	\begin{align*}
		S(t,p)=\sum_{n=1}^{\infty}\frac{e^{\mathrm{i}nt}}{n^p},
		\qquad t\notin2\pi\mathbb Z.
	\end{align*}

	\begin{enumerate}
		\item \textbf{固定 $t\notin2\pi\mathbb Z$,讨论 $p$ 的取值.}

		当 $p>1$ 时,
		\begin{align*}
			\left|\frac{e^{\mathrm{i}nt}}{n^p}\right|=\frac1{n^p},
		\end{align*}
		而 $p$ 级数收敛,故原级数绝对收敛.

		当 $0<p\leq1$ 时,一方面,
		\begin{align*}
			\sum_{n=1}^N e^{\mathrm{i}nt}
			=
			e^{\mathrm{i}t}\frac{1-e^{\mathrm{i}Nt}}{1-e^{\mathrm{i}t}},
		\end{align*}
		故
		\begin{align*}
			\left|\sum_{n=1}^N e^{\mathrm{i}nt}\right|
			\leq \frac{2}{|1-e^{\mathrm{i}t}|}.
		\end{align*}
		又 $n^{-p}$ 单调且趋于 $0$.由 Dirichlet 判别法,原级数收敛.另一方面,绝对值级数就是 $\sum n^{-p}$,故在 $0<p\leq1$ 时发散.因此原级数条件收敛.

		当 $p\leq0$ 时,通项 $e^{\mathrm{i}nt}n^{-p}$ 不趋于 $0$,故原级数发散.

		综上,当 $t\notin2\pi\mathbb Z$ 时,
		\begin{gather*}
			\sum_{n=1}^{\infty}\frac{e^{\mathrm{i}nt}}{n^p}
			\begin{cases}
				\text{绝对收敛},&p>1,\\[3pt]
				\text{条件收敛},&0<p\leq1,\\[3pt]
				\text{发散},&p\leq0.
			\end{cases}
		\end{gather*}

		\item \textbf{固定 $p$,把 $t$ 看成参数.}

		当 $p>1$ 时,
		\begin{align*}
			\left|\frac{e^{\mathrm{i}nt}}{n^p}\right|\leq \frac1{n^p},
		\end{align*}
		且 $\sum n^{-p}$ 收敛.由 Weierstrass 判别法, $S(t,p)$ 在任意 $E\subset\mathbb R$ 上一致绝对收敛.

		当 $0<p\leq1$ 时,对每个固定的 $t\in E$ 级数均收敛,故在 $E$ 上逐点收敛.若 $\delta_E>0$,则
		\begin{align*}
			\left|\sum_{n=1}^N e^{\mathrm{i}nt}\right|
			\leq \frac{2}{|1-e^{\mathrm{i}t}|}
			\leq \frac{2}{\inf_{t\in E}|1-e^{\mathrm{i}t}|},
		\end{align*}
		右端为有限常数.由含参 Dirichlet 判别法, $S(t,p)$ 在 $E$ 上一致收敛.又每个固定 $t$ 对应的绝对值级数均发散,故为一致条件收敛.

		若 $\delta_E=0$,则可取 $t_j\in E$ 与整数 $k_j$ 使
		\begin{align*}
			\eta_j=|t_j-2\pi k_j|\to0.
		\end{align*}
		取 $N_j$ 使 $N_j\eta_j\to1/4$.于是当 $j$ 充分大时,对 $N_j\leq n\leq2N_j$ 有 $n\eta_j\in[1/5,3/5]$,从而这些项的实部同号且有统一正下界.于是尾和
		\begin{align*}
			\sum_{n=N_j}^{2N_j}\frac{e^{\mathrm{i}nt_j}}{n^p}
		\end{align*}
		的模不趋于 $0$: 当 $0<p<1$ 时甚至无界增大,当 $p=1$ 时仍有正下界.由 Cauchy 一致收敛准则, $S(t,p)$ 在 $E$ 上不一致收敛.

		当 $p\leq0$ 时,对每一个 $t\in E$ 通项均不趋于 $0$,故逐点收敛性不成立.

		因此,作为关于 $t$ 的含参级数,
		\begin{gather*}
		\begin{array}{c|c}
			p\text{ 的范围}&S(t,p)\text{ 在 }E\text{ 上的结论}\\ \hline
			p>1&\text{一致绝对收敛}\\[3pt]
			0<p\leq1,\ \delta_E>0&\text{一致条件收敛}\\[3pt]
			0<p\leq1,\ \delta_E=0&\text{逐点收敛但不一致收敛}\\[3pt]
			p\leq0&\text{逐点收敛性不成立}
		\end{array}
		\end{gather*}
	\end{enumerate}
\end{exbox}

\section{无穷积分}

从级数转到积分时,只需把离散尾和换成连续尾积分.例如
\begin{align*}
	\sum_{n=N_1}^{N_2}u_n(t)
	\quad\hbox{对应}\quad
	\int_{A_1}^{A_2}f(x,t)\,\dd x.
\end{align*}
无穷远点的处理完全平行;瑕点也可以通过换元 $x=1/y$ 化为无穷远点,因为
\begin{align*}
	\int_0^a f(x,t)\,\dd x
	=
	\int_{1/a}^{\infty} f\left(\frac1y,t\right)\frac{\dd y}{y^2}.
\end{align*}
所以瑕积分并不是新的收敛机制,只是同一个尾部问题在变量替换后的表现.

\section{基本概念}

\textbf{不含参无穷积分的敛散性.}
设 $f(x)$ 在任意有限区间 $[a,A]$ 上可积.若极限
\begin{align*}
	\lim_{A\to\infty}\int_a^A f(x)\,\dd x
\end{align*}
存在,则称无穷积分
\begin{align*}
	\int_a^\infty f(x)\,\dd x
\end{align*}
收敛;否则称其发散.

若
\begin{align*}
	\int_a^\infty |f(x)|\,\dd x
\end{align*}
收敛,则称原积分绝对收敛;若原积分收敛而绝对值积分发散,则称原积分条件收敛.

\textbf{含参无穷积分的敛散性.}
设 $E$ 为参数集合,且 $f(x,t)$ 对任意 $A>a$ 和任意 $t\in E$ 在 $[a,A]$ 上关于 $x$ 可积.考虑含参无穷积分
\begin{align*}
	\int_a^\infty f(x,t)\,\dd x,
	\qquad t\in E.
\end{align*}
若对每一个固定的 $t\in E$,普通无穷积分都收敛,则称该含参无穷积分在 $E$ 上逐点收敛.

进一步,若对任意 $\varepsilon>0$,存在 $A_0\geq a$,使得当 $A_2>A_1\geq A_0$ 时,对一切 $t\in E$ 均有
\begin{align*}
	\left|\int_{A_1}^{A_2} f(x,t)\,\dd x\right|<\varepsilon,
\end{align*}
则称该含参无穷积分在 $E$ 上一致收敛.这就是含参无穷积分的 Cauchy 一致收敛准则.

若
\begin{align*}
	\int_a^\infty |f(x,t)|\,\dd x
\end{align*}
在 $E$ 上一致收敛,则称原积分在 $E$ 上一致绝对收敛.若原积分在 $E$ 上一致收敛,但绝对值积分在 $E$ 上不一致收敛,则称原积分在 $E$ 上一致条件收敛.特别地,若对每个固定的 $t\in E$,绝对值积分都发散,而原积分在 $E$ 上一致收敛,则原积分必为一致条件收敛.

\textbf{不含参瑕积分的敛散性.}
设 $0<a<\infty$,且 $f(x)$ 在任意闭区间 $[\varepsilon,a]$ 上可积,其中 $0<\varepsilon<a$.若极限
\begin{align*}
	\lim_{\varepsilon\to0^+}\int_\varepsilon^a f(x)\,\dd x
\end{align*}
存在,则称瑕积分
\begin{align*}
	\int_0^a f(x)\,\dd x
\end{align*}
收敛;否则称其发散.绝对收敛与条件收敛的定义同无穷积分.

\textbf{含参瑕积分的敛散性.}
设 $E$ 为参数集合,且 $f(x,t)$ 对任意 $0<\varepsilon<a$ 和任意 $t\in E$ 在 $[\varepsilon,a]$ 上关于 $x$ 可积.考虑含参瑕积分
\begin{align*}
	\int_0^a f(x,t)\,\dd x,
	\qquad t\in E.
\end{align*}
若对每一个固定的 $t\in E$,普通瑕积分都收敛,则称该含参瑕积分在 $E$ 上逐点收敛.

进一步,若对任意 $\varepsilon>0$,存在 $\delta>0$,使得当 $0<\alpha<\beta<\delta$ 时,对一切 $t\in E$ 均有
\begin{align*}
	\left|\int_\alpha^\beta f(x,t)\,\dd x\right|<\varepsilon,
\end{align*}
则称该含参瑕积分在 $E$ 上一致收敛.这就是含参瑕积分的 Cauchy 一致收敛准则.

若
\begin{align*}
	\int_0^a |f(x,t)|\,\dd x
\end{align*}
在 $E$ 上一致收敛,则称原积分在 $E$ 上一致绝对收敛;一致条件收敛的定义与含参无穷积分相同.

\textbf{含瑕点与无穷远点积分的处理.}
对于同时含有瑕点 $0$ 和无穷远点 $\infty$ 的积分
\begin{align*}
	\int_0^\infty f(x)\,\dd x,
\end{align*}
任取 $c>0$,并作分解
\begin{align*}
	\int_0^\infty f(x)\,\dd x
	=
	\int_0^c f(x)\,\dd x+
	\int_c^\infty f(x)\,\dd x.
\end{align*}
若右端两个广义积分均收敛,则称原积分收敛.这一要求与 $c$ 的选取无关.若右端两个绝对值积分均收敛,则称原积分绝对收敛;若原积分收敛而绝对值积分发散,则称原积分条件收敛.

含参情形也按端点分别处理.即
\begin{align*}
	\int_0^\infty f(x,t)\,\dd x
\end{align*}
在 $E$ 上一致收敛,当且仅当对某个,也等价地对任意 $c>0$,
\begin{align*}
	\int_0^c f(x,t)\,\dd x
\end{align*}
在 $0$ 端关于 $t\in E$ 一致收敛,且
\begin{align*}
	\int_c^\infty f(x,t)\,\dd x
\end{align*}
在 $\infty$ 端关于 $t\in E$ 一致收敛.一致绝对收敛也满足相同的分解原则.因此,若某一个端点已经违反相应的 Cauchy 一致收敛准则,则全区间积分不可能在 $E$ 上一致收敛.

\section{常用判别法}

\textbf{比较判别法.}
设当 $x\geq a$ 时 $0\leq f(x)\leq g(x)$,且二者在任意有限区间 $[a,b]$ 上可积.若
\begin{align*}
	\int_a^\infty g(x)\,\dd x
\end{align*}
收敛,则
\begin{align*}
	\int_a^\infty f(x)\,\dd x
\end{align*}
也收敛;若
\begin{align*}
	\int_a^\infty f(x)\,\dd x
\end{align*}
发散,则
\begin{align*}
	\int_a^\infty g(x)\,\dd x
\end{align*}
也发散.

极限形式如下.设当 $x\geq a$ 时 $f(x)\geq0$, $g(x)>0$,且二者在任意有限区间 $[a,b]$ 上可积.若
\begin{align*}
	\lim_{x\to\infty}\frac{f(x)}{g(x)}=L,
\end{align*}
则当 $0<L<\infty$ 时,两积分同时收敛或同时发散.对于 $0$ 到 $a$ 的瑕积分,只需将极限改为 $x\to0^+$.

含参形式如下.若对一切 $x\geq a$ 和一切 $t\in E$ 有
\begin{align*}
	|f(x,t)|\leq \varphi(x),
\end{align*}
且
\begin{align*}
	\int_a^\infty \varphi(x)\,\dd x
\end{align*}
收敛,则
\begin{align*}
	\int_a^\infty f(x,t)\,\dd x
\end{align*}
在 $E$ 上一致绝对收敛.瑕积分的含参形式完全类似,只需把控制函数换成在瑕点附近可积的函数.

\textbf{Dirichlet 判别法.}
设 $f(x),g(x)$ 在 $[a,\infty)$ 上定义,并在任意有限区间上可积.若存在常数 $M>0$,使得对一切 $A\geq a$ 都有
\begin{align*}
	\left|\int_a^A f(x)\,\dd x\right|\leq M,
\end{align*}
又若 $g(x)$ 在 $[a,\infty)$ 上单调,且
\begin{align*}
	\lim_{x\to\infty}g(x)=0,
\end{align*}
则
\begin{align*}
	\int_a^\infty f(x)g(x)\,\dd x
\end{align*}
收敛.

含参形式如下.若存在常数 $M>0$,使得对一切 $A\geq a$ 和一切 $t\in E$ 都有
\begin{align*}
	\left|\int_a^A f(x,t)\,\dd x\right|\leq M,
\end{align*}
并且 $g(x,t)$ 对每个固定的 $t$ 关于 $x$ 单调,且
\begin{align*}
	\lim_{x\to\infty}\sup_{t\in E}|g(x,t)|=0,
\end{align*}
则
\begin{align*}
	\int_a^\infty f(x,t)g(x,t)\,\dd x
\end{align*}
在 $E$ 上一致收敛.

\textbf{Abel 判别法.}
设 $f(x),g(x)$ 在 $[a,\infty)$ 上定义,并在任意有限区间上可积.若
\begin{align*}
	\int_a^\infty f(x)\,\dd x
\end{align*}
收敛,且 $g(x)$ 在 $[a,\infty)$ 上单调有界,则
\begin{align*}
	\int_a^\infty f(x)g(x)\,\dd x
\end{align*}
收敛.

含参形式如下.若
\begin{align*}
	\int_a^\infty f(x,t)\,\dd x
\end{align*}
在 $E$ 上一致收敛,且 $g(x,t)$ 对每个固定的 $t$ 关于 $x$ 单调,并且存在常数 $M>0$,使得
\begin{align*}
	|g(x,t)|\leq M,
	\qquad x\geq a,
	\quad t\in E,
\end{align*}
则
\begin{align*}
	\int_a^\infty f(x,t)g(x,t)\,\dd x
\end{align*}
在 $E$ 上一致收敛.

\textbf{常用幂函数积分.}
设 $0<a<\infty$.则
\begin{align*}
	\int_0^a x^{-p}\,\dd x
\end{align*}
在 $p<1$ 时收敛,在 $p\geq1$ 时发散;而
\begin{align*}
	\int_a^\infty x^{-p}\,\dd x
\end{align*}
在 $p>1$ 时收敛,在 $p\leq1$ 时发散.

\textbf{下面考虑一种典型的积分形式，涵盖上面讨论的所有情况.}

以下固定 $a>0$,并始终假设 $t\neq0$.作为含参积分讨论时,令
\begin{align*}
	E\subset\mathbb R\setminus\{0\},
\end{align*}
并约定
\begin{align*}
	m_E=\inf_{t\in E}|t|,
	\qquad
	M_E=\sup_{t\in E}|t|.
\end{align*}
其中 $m_E=0$ 表示参数可以趋近于 $0$,而 $M_E=\infty$ 表示参数可以无界增大.

此外,后文反复使用常数
\begin{align*}
	C_p=\int_1^2\sin u\,u^{-p}\,\dd u.
\end{align*}
由于 $1<u<2<\pi$,故 $\sin u>0$,从而 $C_p>0$.每当需要证明不一致收敛时,我们都将通过选取小区间 $[1/|t|,2/|t|]$ 并换元 $u=|t|x$ 来得到一个不趋于 $0$ 的 Cauchy 差.

\begin{exbox}
	考虑无穷远端的积分
	\begin{align*}
		I_a(t,p)=\int_a^\infty \sin(tx)x^{-p}\,\dd x.
	\end{align*}

	\begin{enumerate}
		\item \textbf{固定 $t\neq0$,讨论 $p$ 的取值.}

		\begin{enumerate}
			\item 当 $p>1$ 时,
			\begin{align*}
				|\sin(tx)x^{-p}|\leq x^{-p},
			\end{align*}
			而 $\int_a^\infty x^{-p}\,\dd x$ 收敛.由比较判别法,原积分绝对收敛.

			\item 当 $0<p\leq1$ 时,一方面,对任意 $A\geq a$,
			\begin{align*}
				\int_a^A \sin(tx)\,\dd x
				=
				\frac{\cos(ta)-\cos(tA)}{t},
			\end{align*}
			故
			\begin{align*}
				\left|\int_a^A \sin(tx)\,\dd x\right|
				\leq \frac{2}{|t|}.
			\end{align*}
			又 $x^{-p}$ 在 $[a,\infty)$ 上单调且趋于 $0$.由 Dirichlet 判别法,原积分收敛.

			另一方面,
			\begin{align*}
				|\sin(tx)|x^{-p}
				\geq \sin^2(tx)x^{-p}
				=\frac12x^{-p}-\frac12\cos(2tx)x^{-p}.
			\end{align*}
			其中 $\int_a^\infty x^{-p}\,\dd x$ 在 $0<p\leq1$ 时发散,而 $\int_a^\infty \cos(2tx)x^{-p}\,\dd x$ 由同样的 Dirichlet 判别法收敛.故 $\int_a^\infty |\sin(tx)|x^{-p}\,\dd x$ 发散.因此原积分条件收敛.

			\item 当 $p=0$ 时,
			\begin{align*}
				\int_a^A \sin(tx)\,\dd x
				=
				\frac{\cos(ta)-\cos(tA)}{t}.
			\end{align*}
			当 $A\to\infty$ 时右端没有极限,故原积分发散.

			\item 当 $p<0$ 时,反设 $\int_a^\infty \sin(tx)x^{-p}\,\dd x$ 收敛.由于 $x^p$ 在 $[a,\infty)$ 上单调有界,由 Abel 判别法,
			\begin{align*}
				\int_a^\infty \sin(tx)x^{-p}\cdot x^p\,\dd x
				=
				\int_a^\infty \sin(tx)\,\dd x
			\end{align*}
			应当收敛,与 $p=0$ 的结论矛盾.故原积分发散.
		\end{enumerate}

		综上,当 $t\neq0$ 时,
		\begin{gather*}
			\int_a^\infty \sin(tx)x^{-p}\,\dd x
			\begin{cases}
				\text{绝对收敛},& p>1,\\[3pt]
				\text{条件收敛},& 0<p\leq1,\\[3pt]
				\text{发散},& p\leq0.
			\end{cases}
		\end{gather*}

		\item \textbf{固定 $p$,把 $t$ 看成参数.}

		\begin{enumerate}
			\item 当 $p>1$ 时,
			\begin{align*}
				|\sin(tx)x^{-p}|\leq x^{-p},
			\end{align*}
			且 $\int_a^\infty x^{-p}\,\dd x$ 收敛.由一致比较判别法, $I_a(t,p)$ 在任意 $E\subset\mathbb R\setminus\{0\}$ 上一致绝对收敛.

			\item 当 $0<p\leq1$ 时,对每个固定的 $t\in E$,积分均收敛,故在 $E$ 上逐点收敛.

			若 $m_E>0$,则
			\begin{align*}
				\left|\int_a^A\sin(tx)\,\dd x\right|
				\leq \frac{2}{|t|}
				\leq \frac{2}{m_E},
				\qquad t\in E.
			\end{align*}
			又 $x^{-p}$ 单调趋于 $0$.由含参 Dirichlet 判别法, $I_a(t,p)$ 在 $E$ 上一致收敛.由于每个固定的 $t\in E$ 对应的绝对值积分均发散,故为一致条件收敛.

			若 $m_E=0$,则取 $t_n\in E$ 使 $|t_n|\to0$.令
			\begin{align*}
				A_n=\frac1{|t_n|},
				\qquad
				B_n=\frac2{|t_n|}.
			\end{align*}
			于是 $A_n,B_n\to\infty$,且当 $n$ 充分大时 $B_n>A_n\geq a$.由换元 $u=|t_n|x$ 得
			\begin{align*}
				\int_{A_n}^{B_n}\sin(t_nx)x^{-p}\,\dd x
				=
				\operatorname{sgn}(t_n)|t_n|^{p-1}C_p.
			\end{align*}
			因为 $C_p>0$,所以当 $0<p<1$ 时上述尾积分绝对值趋于 $+\infty$,当 $p=1$ 时其绝对值恒为 $C_1>0$.这说明尾部 Cauchy 差不能关于 $t\in E$ 一致趋于 $0$.由 Cauchy 一致收敛准则, $I_a(t,p)$ 在 $E$ 上不一致收敛.

			\item 当 $p\leq0$ 时,对每一个 $t\in E$,积分均发散,故逐点收敛性不成立.
		\end{enumerate}

		因此,作为关于 $t$ 的含参无穷积分,
		\begin{gather*}
		\begin{array}{c|c}
			p\text{ 的范围}&I_a(t,p)\text{ 在 }E\text{ 上的结论}\\ \hline
			p>1&\text{一致绝对收敛}\\[3pt]
			0<p\leq1,\ m_E>0&\text{一致条件收敛}\\[3pt]
			0<p\leq1,\ m_E=0&\text{逐点收敛但不一致收敛}\\[3pt]
			p\leq0&\text{逐点收敛性不成立}
		\end{array}
		\end{gather*}
	\end{enumerate}
\end{exbox}

\begin{exbox}
	考虑 $0$ 端的瑕积分
	\begin{align*}
		J_a(t,p)=\int_0^a \sin(tx)x^{-p}\,\dd x.
	\end{align*}

	\begin{enumerate}
		\item \textbf{固定 $t\neq0$,讨论 $p$ 的取值.}

		由于
		\begin{align*}
			\lim_{x\to0^+}\frac{|\sin(tx)|x^{-p}}{x^{1-p}}=|t|,
		\end{align*}
		且 $0<|t|<\infty$,由比较判别法的极限形式,
		\begin{align*}
			\int_0^a |\sin(tx)|x^{-p}\,\dd x
		\end{align*}
		与
		\begin{align*}
			\int_0^a x^{1-p}\,\dd x
		\end{align*}
		同时收敛或同时发散.而 $\int_0^a x^{1-p}\,\dd x$ 在 $p<2$ 时收敛,在 $p\geq2$ 时发散.

		\begin{enumerate}
			\item 当 $p<2$ 时,绝对值积分收敛,故原积分绝对收敛.

			\item 当 $p\geq2$ 时,取充分小的 $\delta>0$,使 $0<\delta<a$ 且 $\operatorname{sgn}(t)\sin(tx)\geq0$ 在 $0<x<\delta$ 上成立.又
			\begin{align*}
				\lim_{x\to0^+}
				\frac{\operatorname{sgn}(t)\sin(tx)x^{-p}}{x^{1-p}}
				=|t|.
			\end{align*}
			由比较判别法的极限形式,
			\begin{align*}
				\int_0^\delta \operatorname{sgn}(t)\sin(tx)x^{-p}\,\dd x
			\end{align*}
			发散,故 $J_a(t,p)$ 发散.
		\end{enumerate}

		综上,当 $t\neq0$ 时,
		\begin{gather*}
			\int_0^a \sin(tx)x^{-p}\,\dd x
			\begin{cases}
				\text{绝对收敛},&p<2,\\[3pt]
				\text{发散},&p\geq2.
			\end{cases}
		\end{gather*}
		这里没有条件收敛情形,因为在 $0$ 的充分小右邻域内, $\sin(tx)$ 不发生无限振荡抵消.

		\item \textbf{固定 $p$,把 $t$ 看成参数.}

		\begin{enumerate}
			\item 当 $p<1$ 时,对任意 $t\in E$,
			\begin{align*}
				|\sin(tx)x^{-p}|\leq x^{-p},
			\end{align*}
			且 $\int_0^a x^{-p}\,\dd x$ 收敛.由一致比较判别法, $J_a(t,p)$ 在任意 $E\subset\mathbb R\setminus\{0\}$ 上一致绝对收敛.

			\item 当 $1\leq p<2$ 时,对每个固定的 $t\in E$,积分均绝对收敛,故在 $E$ 上逐点收敛.

			若 $M_E<\infty$,则对任意 $t\in E$,
			\begin{align*}
				|\sin(tx)x^{-p}|
				\leq |t|x^{1-p}
				\leq M_E x^{1-p}.
			\end{align*}
			又 $\int_0^a x^{1-p}\,\dd x$ 收敛.由一致比较判别法, $J_a(t,p)$ 在 $E$ 上一致绝对收敛.

			若 $M_E=\infty$,则取 $t_n\in E$ 使 $|t_n|\to\infty$.令
			\begin{align*}
				\alpha_n=\frac1{|t_n|},
				\qquad
				\beta_n=\frac2{|t_n|}.
			\end{align*}
			于是 $0<\alpha_n<\beta_n\to0$,且当 $n$ 充分大时 $\beta_n<a$.由换元 $u=|t_n|x$ 得
			\begin{align*}
				\int_{\alpha_n}^{\beta_n}\sin(t_nx)x^{-p}\,\dd x
				=
				\operatorname{sgn}(t_n)|t_n|^{p-1}C_p.
			\end{align*}
			因为 $C_p>0$,所以当 $p=1$ 时上述积分绝对值恒为 $C_1>0$,当 $1<p<2$ 时其绝对值趋于 $+\infty$.这说明瑕点处的 Cauchy 差不能关于 $t\in E$ 一致趋于 $0$.由 Cauchy 一致收敛准则, $J_a(t,p)$ 在 $E$ 上不一致收敛.

			\item 当 $p\geq2$ 时,对每一个 $t\in E$,积分均发散,故逐点收敛性不成立.
		\end{enumerate}

		因此,作为关于 $t$ 的含参瑕积分,
		\begin{gather*}
		\begin{array}{c|c}
			p\text{ 的范围}&J_a(t,p)\text{ 在 }E\text{ 上的结论}\\ \hline
			p<1&\text{一致绝对收敛}\\[3pt]
			1\leq p<2,\ M_E<\infty&\text{一致绝对收敛}\\[3pt]
			1\leq p<2,\ M_E=\infty&\text{逐点收敛但不一致收敛}\\[3pt]
			p\geq2&\text{逐点收敛性不成立}
		\end{array}
		\end{gather*}
	\end{enumerate}
\end{exbox}

\begin{exbox}
	考虑同时含有瑕点和无穷远点的积分
	\begin{align*}
		K(t,p)=\int_0^\infty \sin(tx)x^{-p}\,\dd x.
	\end{align*}
	任取 $a>0$,拆成
	\begin{align*}
		K(t,p)
		=
		\int_0^a \sin(tx)x^{-p}\,\dd x
		+
		\int_a^\infty \sin(tx)x^{-p}\,\dd x.
	\end{align*}
	因此 $0$ 端由前一个关于瑕点的例子控制, $\infty$ 端由前一个关于无穷远端的例子控制.

	\begin{enumerate}
		\item \textbf{固定 $t\neq0$,讨论 $p$ 的取值.}

		$0$ 端收敛要求 $p<2$,而 $\infty$ 端收敛要求 $p>0$.故原积分收敛当且仅当
		\begin{align*}
			0<p<2.
		\end{align*}
		进一步, $0$ 端只要收敛便绝对收敛,即要求 $p<2$; $\infty$ 端绝对收敛要求 $p>1$.故原积分绝对收敛当且仅当
		\begin{align*}
			1<p<2.
		\end{align*}
		综上,当 $t\neq0$ 时,
		\begin{gather*}
			\int_0^\infty \sin(tx)x^{-p}\,\dd x
			\begin{cases}
				\text{绝对收敛},&1<p<2,\\[3pt]
				\text{条件收敛},&0<p\leq1,\\[3pt]
				\text{发散},&p\leq0\text{ 或 }p\geq2.
			\end{cases}
		\end{gather*}

		\item \textbf{固定 $p$,把 $t$ 看成参数.}

		由含瑕点与无穷远点积分的处理原则, $K(t,p)$ 在 $E$ 上一致收敛,必须同时要求 $0$ 端和 $\infty$ 端都一致收敛;若某一个端点不一致收敛,则整体不一致收敛.

		\begin{enumerate}
			\item 当 $1<p<2$ 时, $\infty$ 端由一致比较判别法在任意 $E$ 上一致绝对收敛; $0$ 端在 $M_E<\infty$ 时由一致比较判别法一致绝对收敛.故当 $M_E<\infty$ 时, $K(t,p)$ 在 $E$ 上一致绝对收敛.

			若 $M_E=\infty$,则 $0$ 端已经不一致收敛.事实上,前一个关于瑕点的例子中构造的区间 $[\alpha_n,\beta_n]\subset(0,a)$ 使得
			\begin{align*}
				\left|\int_{\alpha_n}^{\beta_n}\sin(t_nx)x^{-p}\,\dd x\right|
			\end{align*}
			不趋于 $0$,从而整体积分也违反 $0$ 端的 Cauchy 一致收敛准则.故 $K(t,p)$ 在 $E$ 上不一致收敛.

			\item 当 $0<p<1$ 时, $0$ 端由一致比较判别法在任意 $E$ 上一致绝对收敛; $\infty$ 端在 $m_E>0$ 时由含参 Dirichlet 判别法一致收敛.故当 $m_E>0$ 时, $K(t,p)$ 在 $E$ 上一致收敛.又此时对每个固定的 $t\in E$, $\infty$ 端的绝对值积分发散,所以 $K(t,p)$ 在 $E$ 上一致条件收敛.

			若 $m_E=0$,则 $\infty$ 端已经不一致收敛.由前一个关于无穷远端的例子中构造的尾区间 $[A_n,B_n]$ 可知尾部 Cauchy 差不趋于 $0$,故 $K(t,p)$ 在 $E$ 上不一致收敛.

			\item 当 $p=1$ 时, $0$ 端的一致收敛要求 $M_E<\infty$,而 $\infty$ 端的一致收敛要求 $m_E>0$.故当
			\begin{align*}
				0<m_E\leq M_E<\infty
			\end{align*}
			时, $K(t,1)$ 在 $E$ 上一致收敛.又对每个固定的 $t\in E$, $\infty$ 端的绝对值积分发散,所以 $K(t,1)$ 在 $E$ 上一致条件收敛.

			若 $m_E=0$,则 $\infty$ 端不一致收敛;若 $M_E=\infty$,则 $0$ 端不一致收敛.因此只要 $m_E=0$ 或 $M_E=\infty$, $K(t,1)$ 就在 $E$ 上不一致收敛.

			\item 当 $p\leq0$ 或 $p\geq2$ 时,对每一个 $t\in E$,积分均发散,故逐点收敛性不成立.
		\end{enumerate}

		因此,作为关于 $t$ 的含参广义积分,
		\begin{gather*}
		\begin{array}{c|c}
			p\text{ 的范围}&K(t,p)\text{ 在 }E\text{ 上的结论}\\ \hline
			1<p<2,\ M_E<\infty&\text{一致绝对收敛}\\[3pt]
			1<p<2,\ M_E=\infty&\text{逐点收敛但不一致收敛}\\[3pt]
			0<p<1,\ m_E>0&\text{一致条件收敛}\\[3pt]
			0<p<1,\ m_E=0&\text{逐点收敛但不一致收敛}\\[3pt]
			p=1,\ 0<m_E\leq M_E<\infty&\text{一致条件收敛}\\[3pt]
			p=1,\ m_E=0\text{ 或 }M_E=\infty&\text{逐点收敛但不一致收敛}\\[3pt]
			p\leq0\text{ 或 }p\geq2&\text{逐点收敛性不成立}
		\end{array}
		\end{gather*}
	\end{enumerate}
\end{exbox}

总的来说,对 $\int_0^\infty \sin(tx)x^{-p}\,\dd x$ 而言, $\infty$ 端给出 $p>0$ 的收敛要求和 $p>1$ 的绝对收敛要求; $0$ 端给出 $p<2$ 的收敛要求.作为关于 $t$ 的含参积分时, $t\to0$ 会破坏 $\infty$ 端的统一振荡控制; $|t|\to\infty$ 会破坏 $0$ 端的一致比较控制.

